% ============================================================
%  B 实验操作单（照着做就行）— 大字版，可打印
%  编译：在本目录运行  xelatex "B实验操作单.tex"
%  语言刻意简化，供任何人照着执行
% ============================================================
\documentclass[a4paper, 12pt]{ctexart}
\usepackage{geometry}
\geometry{left=1.6cm, right=1.6cm, top=1.6cm, bottom=1.6cm}
\usepackage{booktabs}
\usepackage{array}
\usepackage{tabularx}
\usepackage{amsmath}
\usepackage{amssymb}
\usepackage{enumitem}
\setlist[itemize]{itemsep=4pt, topsep=2pt}
\setlist[enumerate]{itemsep=4pt, topsep=2pt}
\renewcommand{\arraystretch}{1.45}
\pagestyle{empty}
\newcommand{\ck}{$\square$~}
\begin{document}

\begin{center}
{\LARGE\bfseries 实验操作单（照着做）}\\[6pt]
{\large 样品 B ｜ 用桌上剩下的液体 ｜ 照着做就行（共 4 页）}\\[4pt]
{\normalsize 开始日期：\rule{2.5cm}{0.4pt} \quad 操作人：\rule{3cm}{0.4pt}}
\end{center}

\section*{一、开工前，先把这些放在桌上（放好了打勾）}
\begin{itemize}
  \item \ck 台面上有三个烧杯，贴着标签：\textbf{基液}、\textbf{外层液}、\textbf{中层液}（都是以前配好、还没用完的）
  \item \ck 正丙醇锆（小瓶，怕水，用完马上盖紧）
  \item \ck 葡萄糖（白色粉末，2.00\,g）　\ck 去离子水（3.0\,mL）
  \item \ck 天平（能称到 0.01\,g）、玻璃棒、一次性滴管或注射器
  \item \ck 三个 10\,mL 注射器 + \textbf{全金属针头}（塑料针头会被液体泡坏）
  \item \ck 硅油纸（约 30\,cm 长的一条）、胶带
  \item \ck 静电纺丝机　\ck 马弗炉　\ck 管式炉（通氩气那台）
  \item \ck 笔和本子（下面第 5 节要填）
\end{itemize}

\section*{二、时间表：什么时候干什么}

\begin{table}[h]
\centering
\begin{tabularx}{\textwidth}{p{2.6cm} X}
\toprule
\textbf{时间} & \textbf{要做的事} \\
\midrule
\textbf{第 1 天}\newline 上午 9:00 & \textbf{（1）先检查液体。} 三个烧杯都举起来看一看、用玻璃棒挑一挑：\\ 
 & 　对的样子：透明、不浑、没有白色絮状、没有胶块，挑起来是短短的一小段丝。\\
 & \textbf{如果发浑、有胶块、或者一挑就拉长丝不断，立刻停下}，告诉负责人（这批液体放了 18 天，可能坏了）。\\
\midrule
\textbf{第 1 天}\newline 上午 9:20 & \textbf{（2）试纺。} 倒 2--3\,mL 出来，在小纸片上纺一下，看喷出来的丝稳不稳（不停不断、不滴液）。稳了再往下做。\\
\midrule
\textbf{第 1 天}\newline 上午 9:40 & \textbf{（3）加锆。} 用注射器往\textbf{中层液}里慢慢滴入正丙醇锆 \textbf{0.64\,g}，一边滴一边搅，滴完继续搅 \textbf{30 分钟}。\\
 & 　对的样子：液体还是透明，没有变浑、没有白点。有白点就停下报告。\\
\midrule
\textbf{第 1 天}\newline 上午 10:20 & \textbf{（4）配糖水。} 葡萄糖 2.00\,g + 去离子水 3.0\,mL，放进 40--50\,°C 的温水里隔着温热，搅到完全化开、没有颗粒，放到不烫手。\\
\midrule
\textbf{第 1 天}\newline 上午 10:40 & \textbf{（5）把糖水倒进中层液}，搅 \textbf{15--30 分钟}。\\
 & \textbf{【最重要】搅完 30 分钟以内必须开始纺丝}，不能放着过夜。\\
\midrule
\textbf{第 1 天}\newline 上午 11:10\newline 到第 2 天早上 & \textbf{（6）纺丝，大约 20 小时，中间不能停。} 顺序：\\ 
 & 　先纺\textbf{外层液} 10\,mL，再纺\textbf{中层液} 10\,mL，最后再纺\textbf{外层液} 10\,mL。\\
 & 　出液速度 1.5\,mL/h（机器设定好不要改）。硅油纸接住。\\
 & 　每层开始前先空喷一小会儿，等丝稳定了再让硅油纸接。\\
 & 　一晚上没人也要开着；中途去看一眼，正常就是细细的丝不断。\\
\midrule
\textbf{第 2 天}\newline 早上 7:00 & \textbf{（7）取膜。} 把纺好的膜连着硅油纸拿下来，平放，称一次重量（\textbf{初重}）记下来，晾 1--2 小时。不要折叠。\\
\midrule
\textbf{第 2 天}\newline 下午 14:00 & \textbf{（8）空气预氧化（马弗炉，不盖任何东西）。} 把膜平放进马弗炉：\\
 & 　室温慢慢升到 \textbf{100\,°C}，停 30 分钟；\\
 & 　再升到 \textbf{150\,°C}，停 30 分钟；\\
 & 　再升到 \textbf{220\,°C}，停 60 分钟；\\
 & 　然后关火，等炉子自己凉到 \textbf{60\,°C 以下}再开门取出来。\textbf{整个过程不能超过 220\,°C。}\\
 & 　对的样子：膜的颜色变成\textbf{均匀的金黄色或浅褐色}，还是完整的一片。称一次重量记下（预氧化后重）。\\
\midrule
\textbf{第 2 天}\newline 傍晚 18:00 & \textbf{（9）进管式炉（通氩气）。} 把膜平放进瓷舟（写上编号），放进管式炉，先\textbf{通氩气 30 分钟}（把空气赶走），然后开程序：\\
 & 　200 分钟升到 \textbf{200\,°C}，停 60 分钟；\\
 & 　400 分钟升到 \textbf{900\,°C}，停 60 分钟；\\
 & 　然后关火，让它自己慢慢凉（整夜）。\\
\midrule
\textbf{第 3 天}\newline 早上 & \textbf{（10）出炉。} 等炉子凉到 \textbf{100\,°C 以下}再取。称一次重量（\textbf{终重}）记下来。\\
 & 　把样品装进样品袋，写好标签：\textbf{B + 日期 + 名字}。\\
\midrule
\textbf{第 3 天}\newline 上午 & \textbf{（11）拍照 + 送检。} 拍两张照片（整片的样子、用手弯一下的样子）。联系送 \textbf{SEM}，说明"要看表面，如果能切就再看截面"。\\
\bottomrule
\end{tabularx}
\end{table}

\section*{三、绝对不能做错的 6 件事}
\begin{enumerate}
  \item \textbf{顺序不能反}：先加锆，后加糖水。反过来锆会析出白点。
  \item \textbf{糖水加进去 30 分钟内必须开始纺丝}，不能放。
  \item \textbf{针头必须全金属}，塑料的会被液体泡坏、堵住。
  \item \textbf{预氧化不能超过 220\,°C}，也不能升温太快，否则膜会烧黑、发脆。
  \item \textbf{两个炉子的顺序不能反}：先马弗炉（有空气），后管式炉（通氩气）。换炉子前必须先凉到 60\,°C 以下。
  \item \textbf{管式炉一定要先通氩气 30 分钟}，把里面的空气赶走再升温。
\end{enumerate}

\section*{四、出问题了怎么办（照着做）}
\begin{itemize}
  \item 加锆后\textbf{出现白点、变浑} $\rightarrow$ 马上停下，不要纺，记下来报告。
  \item 纺丝时\textbf{老是断、老是滴} $\rightarrow$ 先检查针头堵没堵，换个针头再试。
  \item 纺出来的膜\textbf{太薄、像没纺上} $\rightarrow$ 叫负责人来看，先别关机器。
  \item 预氧化后\textbf{颜色发黑、一碰就碎} $\rightarrow$ 说明温度高了，记下来，下次降到 200\,°C。
  \item 液体\textbf{变得很稠、一挑就是长丝} $\rightarrow$ 记下时间，不要再硬纺，报告负责人（可能要重新配液体）。
\end{itemize}

\section*{五、要填的数字（做完一项填一项）}
\begin{table}[h]
\centering
\begin{tabularx}{\textwidth}{X p{3.2cm} p{3.2cm}}
\toprule
\textbf{项目} & \textbf{要求} & \textbf{实际} \\
\midrule
三个烧杯的检查结果 & 透明 / 不透明 & \\
试纺是否稳定 & 稳定 / 不稳 & \\
正丙醇锆 称了多少克 & 0.64\,g & \\
葡萄糖 / 水的量 & 2.00\,g / 3.0\,mL & \\
搅完糖水到开始纺丝隔多久 & 30 分钟以内 & \\
纺丝开始时间 / 结束时间 & —— & \\
初重（纺完） & —— & \\
预氧化后的颜色 & 金色 / 褐色 & \\
预氧化后重量 & —— & \\
终重（烧完） & —— & \\
\bottomrule
\end{tabularx}
\end{table}

\vspace{4pt}
\noindent\textbf{做完把这张纸和第 5 节的数字一起交给负责人。}

\end{document}
